\section{Estimand and estimator}
\label{sec:estimator}
\subsection{Runs and score scales}
Let $c$ index configurations, $r$ index their $R$ replicate runs, and $j$ index benchmarks, where a configuration fixes a data recipe, model size, and nominal hyperparameters, and Table~\ref{tab:notation} in Appendix~\ref{app:derivations} collects the symbols of the construction and its diagnostics. The seed in DataDecide controls initialisation and data order together, and auxiliary training schedules introduce an additional source of run variation. For item $i$, define the per-byte margin and corresponding correctness indicator by
\begin{equation}
m_{crji}=\frac{\ell(\mathrm{gold}_i)}{\mathrm{bytes}(\mathrm{gold}_i)}-
\max_{v\ne\mathrm{gold}_i}\frac{\ell(v_i)}{\mathrm{bytes}(v_i)},
\qquad a_{crji}=\mathbf{1}\{m_{crji}>0\}.
\label{eq:scores}
\end{equation}
In Equation~\ref{eq:scores}, $\ell$ is the summed answer log-likelihood and $v$ ranges over alternatives, with a zero margin treated as incorrect. Items are averaged within each benchmark, with MMLU macro-averaged over its 57 subjects, and we denote the resulting margin and accuracy scores by $y^{\mathrm{marg}}$ and $y^{\mathrm{acc}}$, calling each benchmark-level score a trait.

For each half $H\in\{A,B\}$, our measurement model is
\begin{equation}
y^H_{crj}=\mu^H_{cj}+E_{crj}+\varepsilon^H_{crj}.
\end{equation}
A half-specific mean lets the two item sets differ in difficulty, the component $E_{crj}$ is the run deviation shared by the two item populations, and the term $\varepsilon^H_{crj}$ carries the remaining item-dependent variation. When deriving unbiasedness we assume independent, identically distributed run effects within a configuration, and we write their covariance as $\Sigma_{E,c}$. However, we can't sustain that interpretation in full, because training schedules differ within the actual population.

\subsection{Cross-half covariance}
We partition each benchmark's items into disjoint halves with a frozen random seed and reuse that assignment for its replicate runs. Each half score is then centred across the replicate runs, which defines $d^H_{crj}=y^H_{crj}-R^{-1}\sum_s y^H_{csj}$. If the item errors have conditional mean zero and zero covariance between the two halves across the relevant runs and benchmarks, then
\begin{equation}
\mathbb{E}[d^A_{crj}d^B_{crk}]=(1-1/R)\Sigma_{E,c}(j,k).
\label{eq:identity}
\end{equation}
Equation~\ref{eq:identity} covers diagonal entries without requiring Gaussian run effects. However, disjoint item identifiers alone don't establish the required independence, since passages, duplicate content and other shared item effects can connect the halves, a risk we check for the one benchmark with documented shared passages by assigning whole passages to halves in Section~\ref{sec:composition}. Item noise shared across benchmarks and repeated across halves would also enter the numerator, and we therefore read our estimate as identified only up to that component, since Appendix~\ref{app:calibration} shows a 0.124 share of item variance would reproduce 1.244.

We write the average centred score as $g^H_{cr}=K^{-1}\sum_j d^H_{crj}$ and compute
\begin{equation}
T_c=\frac{1}{R-1}\sum_r g^A_{cr}g^B_{cr},\qquad
U_c=\frac{1}{K^2(R-1)}\sum_{j,r}d^A_{crj}d^B_{crj}.
\label{eq:tu}
\end{equation}
The respective expectations of the two statistics in Equation~\ref{eq:tu} are $K^{-2}\mathbf{1}^{\mathsf T}\Sigma_{E,c}\mathbf{1}$ and $K^{-2}\operatorname{tr}\Sigma_{E,c}$. Each statistic is summed across configurations before the division, which gives
\begin{equation}
\widehat\Lam=\sqrt{\frac{\sum_c T_c}{\sum_c U_c}},\qquad
\Lam^2=\frac{\mathbf{1}^{\mathsf T}\Sig\mathbf{1}}{\operatorname{tr}\Sig}
=1+(K-1)\bar r_E,\qquad K_{\mathrm{eff}}=\frac{K}{\Lam^2},
\label{eq:lambda}
\end{equation}
where $\Sig=N^{-1}\sum_c\Sigma_{E,c}$ is the covariance matrix averaged across configurations, whose off-diagonal contribution we express through the effective coefficient
\begin{equation}
\bar r_E=\frac{\sum_{j\ne k}\Sigma_E(j,k)}{(K-1)\operatorname{tr}\Sig}.
\label{eq:effective}
\end{equation}
This effective covariance coefficient equals the mean off-diagonal correlation when the marginal variances are equal, and with unequal variances it differs from a normalised weighted mean correlation (Appendix~\ref{app:derivations}). Note that our $K_{\mathrm{eff}}$ compares this battery against independent benchmarks of the same average variance and doesn't count independent information sources.

Unbiased estimates of $T_c$ and $U_c$ don't make their ratio or square root unbiased, and either pooled sum can be nonpositive in a finite sample, which is why we use calibration simulations to check finite-sample behaviour in specified populations.

\begin{algorithm}[t]
\caption{SNAP measurement recipe for a fixed battery of $K$ benchmarks scored on $R$ replicate runs per configuration. The primary analysis below used a random frozen split and checked a passage-aware split, and the registered test kept passage-linked items in one half.}
\label{alg:snap}
\begin{algorithmic}[1]
\STATE Score every replicate run $r$ of every configuration $c$ on every item of the $K$ benchmarks.
\STATE Assign each benchmark's items to halves $A$ and $B$ once, with a frozen seed, and reuse that assignment for every run.
\STATE Average items within each half to obtain $y^A_{crj}$ and $y^B_{crj}$, then centre across the $R$ runs within each configuration to obtain $d^H_{crj}$.
\STATE Form the battery averages $g^H_{cr}$ and the statistics $T_c$ and $U_c$ of Equation~\ref{eq:tu}.
\STATE Pool over configurations and report $\widehat\Lam=\sqrt{\sum_c T_c/\sum_c U_c}$ and $K_{\mathrm{eff}}=K/\widehat\Lam^2$.
\STATE Report a wild cluster bootstrap-$t$ interval over the recipe clusters (Appendix~\ref{app:implementation}) and repeat the split to measure re-split variation.
\STATE Recompute $\widehat\Lam$ after removing each benchmark and each cluster in turn, and report that spread beside the interval, since the factor belongs to the battery.
\end{algorithmic}
\end{algorithm}


\subsection{Resolution and inference}
For an infinitesimal additive shift in Gaussian item margins, the ratio of accuracy to margin noise per unit squared sensitivity is
\begin{equation}
\frac{p(1-p)}{\phi(z_0)^2},\qquad z_0=\Phi^{-1}(1-p).
\label{eq:information}
\end{equation}
At an illustrative accuracy of $p=0.35$ this model gives about 1.66, and its minimum is $\pi/2$, which means the planned lower threshold of 1.4 for per-trait ratios couldn't falsify the model (prediction eight of the plan, Appendix~\ref{app:design}). Yet the calculation describes only the assumed location-shift model and doesn't show an empirical information advantage for arbitrary margin distributions.

We report nominal 95\% wild cluster bootstrap-$t$ intervals over the 25 recipe clusters, using 4,999 Rademacher draws and a linearisation of the squared ratio. Appendix~\ref{app:implementation} states the implementation and its fixed observed denominator during resampling. We also report a configuration percentile bootstrap that resamples both sums but ignores recipe dependence, and neither interval width alone identifies a design effect.

We report the wild interval because it covers well across the twelve populations that Table~\ref{tab:coverage} in the appendix lists. A 2,000-replicate simulation per population, using the same 4,999 draws, gives wild bootstrap coverage of 0.928 to 0.954, while the configuration percentile bootstrap drops to 0.872 under recipe-shared effects. Cross-half error correlation of 0.1 biases the estimate downward by 0.014 at margin-like noise and 0.016 at accuracy-like noise, and under accuracy-like item noise the null population carries a downward bias of 0.011. Finite-sample bias therefore understates an excess rather than manufacturing one.

Our remaining checks leave both verdicts in place. A cluster test-inversion interval that we added later covers 0.935 to 0.956 in a 10,000-replicate repeat of the same populations, gives \ci{1.095}{1.323} for margins and \ci{0.995}{1.162} for accuracy on the observed data, and moves the margin lower endpoint down by 0.048. A residual-scaled wild interval covers 0.943 to 0.958 there and moves neither interval across one, but because we chose it after seeing four of the cells, we report the shipped interval throughout (Appendix~\ref{app:calibration}). Recipe influence on the margin estimate is concentrated, with a participation ratio of 3.08 across the 25 recipes and one recipe holding 0.520 of the squared influence, and removing that recipe moves the estimate by 0.037. With clusters on the five size bands instead of the 25 recipes, the margin interval runs 1.030 to 1.426 and still excludes one, while the accuracy interval runs 0.961 to 1.184 and still includes one (Appendix~\ref{app:supplementary}).

\section{Data and analysis scope}
DataDecide supplies the ten benchmarks ARC-Challenge, ARC-Easy, BoolQ, CommonsenseQA, HellaSwag, MMLU, OpenBookQA, PIQA, Social IQa, and WinoGrande. We analyse 37,682 items per run across 25 recipes at 150M, 300M, 530M, 750M, and 1B parameters, since smaller released models lack the required per-item replicate coverage. Replicate seed labels are $\{2,14,15\}$ below 1B and $\{2,4,5\}$ at 1B.

In each configuration we select the largest checkpoint step available for all three runs. The release's checkpoint index scores this step at an average of 41.5\% of the 750M default runs' final step (range 18\% to 44\%) and at 87\% for the 530M default runs (range 39\% to 89\%). DataDecide stops the auxiliary replicates below 1B at 25 percent of the 1B compute budget. In 26 configurations, one at 530M and 25 at 750M, the largest final step exceeds the smallest by more than a factor of 1.5. At 750M the default runs stop well short of their released final models, whereas the auxiliary runs sit close to their own final steps in every 750M configuration but one. Truncation in our sample therefore concentrates in one band.

We interpret the full estimate as covariance among the released runs at the selected steps, and we report an auxiliary-run contrast that removes the default-versus-auxiliary batch component. Dropping the 26 severely truncated configurations removes the 750M band, giving margin inflation 1.276 (\ci{1.114}{1.420}) and accuracy inflation 1.120 (\ci{1.020}{1.208}), a subset we chose after seeing the full-sample estimates. On the 33 configurations that share a final step, margin inflation is 1.082 (\ci{0.121}{1.530}) and accuracy 1.129 (\ci{0.769}{1.397}), and both intervals include one. Appendix~\ref{app:schedule} finds no detectable rank correlation between a configuration's influence on the estimate and its truncation. At the adjacent earlier shared step the accuracy interval is \ci{1.044}{1.210} and excludes one (Appendix~\ref{app:checkpoint}), and on the 123 configurations that share an adjacent step, the paired change in accuracy inflation between the two steps is 0.048 (\ci{-0.022}{0.130}), which reaches zero. That check can't see an offset shared by every configuration, and our auxiliary contrast doesn't remove differences shared with the checkpoint schedules. We therefore read our accuracy bound as conditional on the checkpoint we scored.

The analysis plan assigned eight recipes to screening and seventeen to estimation, but we used all recipes without drawing that partition, and every choice in the primary analysis could see all 125 configurations. The plan also has neither a supplied original file nor an independently corroborated date. We find that no estimation set reaches the plan's thresholds of 1.349 and 1.40, since an enumeration over the $\binom{25}{17}=1{,}081{,}575$ possible estimation sets produces margin estimates from 1.037 to 1.317 and accuracy estimates from 0.963 to 1.191. Yet this retrospective calculation can't restore a holdout or undo choices made after viewing the data.

We retain the full ten-benchmark battery for our primary estimand, including benchmarks with negative estimated diagonal covariance, since an unbiased moment estimate can be negative near zero and deleting a benchmark because we saw its negative estimate would change both the selection rule and the estimand. Our standardised diagnostics inherit that problem, because they keep only traits with positive variance.

We add a second model family whose replicate design DataDecide lacks. PolyPythias \citep{vanderwal2025} trains nine seeds at each of 14M, 31M, 70M, 160M, and 410M parameters with the same code, data, and hyperparameters, where the seed sets initialisation and data order together, and we score all 45 runs at the shared final step 143,000 on the full battery of 37,682 items. The request text comes from the DataDecide release itself, so every prompt matches the one the DataDecide checkpoints were scored with byte for byte. Our earlier transfer check scored a nested sample of 4,755 items and stays in Appendix~\ref{app:transport}.

For the identification check we built a second bank of 6,808 items from the same ten benchmarks with OLMES at the commit DataDecide used, taking the train split for eight benchmarks and the validation split for OpenBookQA and MMLU. We dropped any item whose context matches a bank-one context or whose question appears among its own few-shot exemplars, and we capped each benchmark at 600 items apart from MMLU. Both reading rules in Section~\ref{sec:family} were committed to the repository before any of these runs was scored, and like the protocol's, that commit time isn't independently corroborated.

